\subsection{Limitations and Future Works}

\subsubsection{Limited Perch Shapes}

The testing in this thesis only focused on detecting long, cylindrical perches without bends or branches.
Since the Hough transform used to detect possible perches is made to detect such shapes, it is likely to
perform poorly if a different shape of perch is introduced, for example a crooked or Y-shaped branch.   

The testing performed on straight perches only is still believed to be a valid proof of concept. For the 
perching mechanism to attach to a perch the drone has to approach a straight segment of the perch, at which
point the visible portion of the perch will appear as a straight object without bends or branches. 

\subsubsection{Limited Field of View}

The cameras used in this device, as positioned in the proposed design, struggle to detect small distances. 
Even after moving the cameras closer together, the device cannot detect distances smaller than roughly 40mm
without the \texttt{StereoSGBM} block matcher cutting off half of the image. 

Refining the circuit layout of the device to integrate the camera sensors on the same PCB as the CPU,
rather than placing them on separate modules, could alleviate this issue by placing the cameras closer 
together. This could also help reduce weight by reducing the number of PCBs and connectors needed. 

Another way to solve this could be to dynamically reconfigure the \texttt{StereoSGBM} parameters in 
a staged process as the perch gets closer to the cameras. When a perch has not been found, or one is found but far away, 
the fullest field of view is needed to find perches in the largest possible area. But once a perch is identified and the 
drone has moved near it, it is known to be near the center of the frame. In this it may be acceptable to discard
many pixels in the margin, since only the center part of the image is needed to detect precise positioning of
the perch. So the detection software could adjust the range of disparities that are measured to track only a window of depth
around where a perch was last detected. This approach could harm runtime performance due to the cost of reconfiguring the 
stereo matchers, though some gains could also be realized by limiting the number of disparities that need to be measured 
in each update.

\subsubsection{Poor Background Separation}

The separation of background pixels performed in this thesis is based on the assumption that 
perches will always be darker than the background. This may not always hold up in reality, 
and in many scenes background points are incorrectly detected as foreground (this can 
be seen in Fig. \ref{fig:bad-downsample:fixed}). Improving the reliability of this method
across a wider range of backgrounds, or enabling it to work in a more cluttered scene,
requires a better method of background removal.

Machine learning based methods could be used to more precisely separate pixels that are 
background, if sufficient data were available. Alternatively, better means of estimating the 
accuracy of disparity data could be used to reject low confidence pixels (since rejection of 
low-quality data is what the background removal achieved). One method to do this could be
to compare the disparity data against time series data across consecutive image frames, 
such as optical flow. Optical flow, paired with IMU data encoding the drones orientation 
and motion, could provide additional depth estimation data to separate out the highest quality
foreground pixels.

\subsubsection{Poor Separation of Lines from Other Shapes}

The Hough transform used has a critical shortcoming, which is that it does not 
distinguish lines from other shapes; it generally works with the assumption that 
any region containing many points is a line. The check added in this proposed workflow which requires 
a line's ratio of length to width be above a parameterized value was not effective in solving this 
issue. This can be seen in Fig. \ref{fig:filter-compare}
and \ref{fig:bad-downsample}, where in both examples a flat background plane has many 
lines drawn through it by the Hough transform. This would need to be fixed for the proposed 
method to work in more cluttered environments.

One method to solve this could be to add a minimum uniqueness limit to the Hough transform 
result filtering. This would involve checking, for each detected line, how many votes are
lost if its anchor point or direction were shifted by a small step, and rejecting any line 
that does not lose a minimum number of votes. This would work to pass shapes that are truly 
long thin lines, and reject other shapes like planes, boxes, or spheres that would still 
have a high number of votes if the line drawn through them were shifted or turned slightly.
